\chapter{概率收敛}

\begin{solution}{probconv-bounded-distance}
记 $D_n=\|X_n-X\|$。对 $0<\varepsilon<1$，
\[
 \varepsilon\mathbb P(D_n>\varepsilon)
 \le\mathbb E\min(1,D_n)
 \le\varepsilon+\mathbb P(D_n>\varepsilon).
\]
左边给出期望趋零蕴含依概率收敛；右边先令 $n\to\infty$ 再令
$\varepsilon\downarrow0$，得到反向。超过或等于 $1$ 的阈值由更小阈值
的估计控制。

定义 $d(X,Y)=\mathbb E\min(1,\|X-Y\|)$。它对称、非负，与可测代表的
a.s. 修改无关；非负函数积分为零的判别给出 $d(X,Y)=0$ 当且仅当
$X=Y$ a.s.。又对 $a,b\ge0$ 有
$\min(1,a+b)\le\min(1,a)+\min(1,b)$，结合 Euclid 三角不等式并积分，
便得 $d(X,Z)\le d(X,Y)+d(Y,Z)$。所以它在等价类上是度量。
\end{solution}

\begin{solution}{probconv-summable-errors}
对事件 $A_n=\{\|X_n-X\|>a_n\}$，第一 Borel--Cantelli 引理给出
$\mathbb P(A_n\text{ 无穷多次})=0$。因此几乎每个样本点都有一个可随该
样本点变化的起点，使以后误差均不超过 $a_n$；由于 $a_n\to0$，得到
a.s. 收敛。

第二问不必先选同一个变化阈值。对每个整数 $m\ge1$，Markov 不等式给
\[
 \sum_n\mathbb P(\|X_n-X\|>1/m)
 \le m^p\sum_n\mathbb E\|X_n-X\|^p<\infty.
\]
故对每个 $m$，误差超过 $1/m$ 只发生有限次几乎必然。对 $m$ 的例外
零集取可数并，便得到同一个概率为一的集合，其上误差趋于零。
\end{solution}

\begin{solution}{probconv-separate-spaces}
对 $0\le t<1$，有
\[
 F_n(t)=\frac{\lfloor nt\rfloor}{n},\qquad
 |F_n(t)-t|\le\frac1n.
\]
当 $t<0$ 时 $F_n(t)=0$，当 $t\ge1$ 时 $F_n(t)=1$。因此分布函数在
每个实点趋于均匀分布的连续分布函数，实线判别给出依分布收敛。
也可对连续测试函数使用区间上的 Riemann 和。

几乎必然收敛涉及同一 $\omega$ 上的整条数列，需要先指定共同样本空间
以及各 $X_n$ 的联合实现。原题只给出了不同有限空间及各自的边缘律，
没有提供这种对象，所以不能直接形成 $X_n(\omega)-X(\omega)$。
这不妨碍已经证明的分布收敛。
\end{solution}

\begin{solution}{probconv-cdf-atom-boundary}
题设说明 $a,b$ 是 $X$ 的分布函数 $F$ 的连续点。因此
\[
 \mathbb P(a<X_n\le b)=F_n(b)-F_n(a)
 \longrightarrow F(b)-F(a)=\mathbb P(a<X\le b).
\]
若取 $X_n=1/n$、$X=0$、$(a,b]=(0,1]$，则 $X_n\to0$ 甚至处处成立，
但左边概率恒为 $1$，极限事件的概率为 $0$。失败来自左端点的原子。
右端点也可出现同样问题，例如 $(a,b]=(-1,0]$ 配合同一变量列。
\end{solution}

\begin{solution}{probconv-lp-strict-hierarchy}
在 $(0,1)$ 上取 Lebesgue 概率，令
$X_n=n^{1/q}\mathbf1_{(0,1/n)}$。则
\[
 \mathbb E|X_n|^q=1,\qquad
 \mathbb E|X_n|^p=n^{p/q-1}\longrightarrow0,
\]
因为 $p/q<1$。每个变量的所有有限阶矩都有限。对每个 $t>0$，当
$n>1/t$ 时 $X_n(t)=0$，所以它还处处收敛到零。这个构造同时说明
a.s. 收敛与较低阶矩收敛并不迫使更高阶矩收敛。
\end{solution}

\begin{solution}{probconv-mean-square-covariance}
正文的乘积估计给出
$\mathbb E(X_nY_n)\to\mathbb E(XY)$。$L^2$ 收敛又蕴含 $L^1$ 收敛，
故 $\mathbb E X_n\to\mathbb E X$、$\mathbb E Y_n\to\mathbb E Y$。
在协方差公式中相减，得到所需结论。

只依概率收敛时，取 $(0,1)$ 上的
$X_n=Y_n=\sqrt{na_n}\mathbf1_{(0,1/n)}$，而 $X=Y=0$，其中
$a_n=1$ 当 $n$ 为偶数，$a_n=2$ 当 $n$ 为奇数。
两个序列都处处趋零，每项均平方可积；但
\[
 \operatorname{Cov}(X_n,Y_n)=\operatorname{Var}(X_n)
 =a_n\left(1-\frac1n\right).
\]
偶数子列和奇数子列分别趋于1和2，所以协方差没有极限。
独立性并未作为这两问的假设，不能在反例中擅自添加它。
\end{solution}

\begin{solution}{probconv-random-limit-coupling}
取公平符号 $Z$，令 $X=Z$、$X_n=-Z$。各变量的律都是
$(\delta_{-1}+\delta_1)/2$，而 $|X_n-X|=2$ 处处，所以题中概率为一。
若共同分布是 $\delta_c$，则对每个 $n$，$X_n=c$、$X=c$ 都几乎必然
成立。有限并的零集说明 $X_n=X$ a.s.，对可数个 $n$ 再取并甚至得到
整列在一个概率为一的集合上都等于 $c$。所以不可能产生这种失败。

分布极限唯一，是说所有候选极限随机变量有相同的概率分布；同一个
分布可以由不同函数实现，既不保证这些函数相等，也不保证它们在一个
预先给定的空间内彼此接近。
\end{solution}

\begin{solution}{probconv-tightness-moment-bound}
若 $C=\sup_n\mathbb E\|X_n\|^p<\infty$，则 Markov 不等式给出
$\sup_n\mathbb P(\|X_n\|>R)\le C/R^p$。给定误差，选充分大的 $R$，
有限维闭球为紧集，故分布族一致紧。

反向矩界不成立。令 $X_n$ 以概率 $1/n$ 取 $n^2$，其余概率取 $0$。
则 $\mathbb E|X_n|=n$。若 $n^2\le R$，其球外概率为零；否则
$\mathbb P(|X_n|>R)=1/n<1/\sqrt R$，故仍一致紧。
最后，常数变量 $Y_n=0$（$n$ 偶数）和 $Y_n=1$（$n$ 奇数）的分布都
支撑在紧集 $\{0,1\}$ 上，却不整列弱收敛。例如连续有界函数
$h(x)=x/(1+|x|)$ 的期望在 $0$ 与 $1/2$ 间交替。
\end{solution}

\begin{solution}{probconv-portmanteau-boundary}
集合包含 $B^\circ\subseteq B\subseteq\overline B$ 先给出对应每项概率
的不等式。对开集 $B^\circ$ 应用开集判据，对闭集 $\overline B$ 应用
闭集判据，再与 $\liminf\le\limsup$ 合并，即得整条不等式。

取 $X_n=X=0$，$B=\{0\}$。则 $\partial B=\{0\}$ 具有极限概率一，
但各集合概率始终为一，当然收敛。因此边界零测保证收敛，却不是每个
具体序列中收敛的必要条件。
\end{solution}

\begin{solution}{probconv-discontinuous-map}
函数 $g(x)=\mathbf1_{(-\infty,t]}(x)$ 仅在 $t$ 处不连续。
连续映射定理的依概率版本给出 $g(X_n)\xrightarrow{\mathbb P}g(X)$。
由于两个示性函数的差只取 $-1,0,1$，对任意 $p>0$，
\[
 \mathbb E|g(X_n)-g(X)|^p
 =\mathbb P(g(X_n)\ne g(X))
 =\mathbb P(|g(X_n)-g(X)|>1/2)\longrightarrow0.
\]
这证明每个有限 $L^p$ 的收敛。若 $X_n=1/n$、$X=t=0$，则 $X_n\to X$
处处，但 $g(X_n)=0$、$g(X)=1$，说明无原子条件不可直接删去。
\end{solution}

\begin{solution}{probconv-joint-marginals}
因为公平符号 $Z$ 与 $-Z$ 同分布，$X_n=Z$ 和 $Y_n=(-1)^nZ$ 的边缘
分布均不随 $n$ 改变。然而对有界连续函数
$h(x,y)=xy/(1+|xy|)$，有
\[
 \mathbb E h(X_n,Y_n)=\frac{(-1)^n}{2},
\]
因此联合分布不能弱收敛。

若两个分量分别依概率收敛，则由
\[
 \|(X_n-X,Y_n-Y)\|\le\|X_n-X\|+\|Y_n-Y\|
\]
及并集上界，联合误差超过 $\varepsilon$ 的概率不大于两个分量误差
超过 $\varepsilon/2$ 的概率之和，故趋零。该论证也适用于有限维向量块。
\end{solution}

\begin{solution}{probconv-independent-joint-limits}
在 $(\mathbb R^2,\mathcal B(\mathbb R^2),\mu\times\nu)$ 上取坐标变量
$X,Y$，则它们独立且有指定分布。记 $X_n,Y_n$ 的分布为 $\mu_n,\nu_n$。
它们的联合律为 $\mu_n\times\nu_n$。

固定 $h\in C_b(\mathbb R^2)$ 和 $\varepsilon>0$。由边缘分布收敛的一致
紧性，可取 $a<b$、$c<d$，使各分布在相应区间外的概率均小于
$\varepsilon$，并且 $a,b$ 是 $\mu$ 的分布函数连续点，$c,d$ 是 $\nu$
的分布函数连续点。这可先取足够大的紧区间，再稍微向外移动端点避开
至多可数个原子。因此各联合律在 $K=(a,b]\times(c,d]$ 外的概率小于
$2\varepsilon$，极限律也如此。

在闭矩形 $[a,b]\times[c,d]$ 上，$h$ 一致连续。取足够细的有限矩形
网格，使每格上振幅小于 $\varepsilon$，且各竖直端点避开 $\mu$ 的
原子、各水平端点避开 $\nu$ 的原子。每个格子 $I\times J$ 的质量满足
\[
 (\mu_n\times\nu_n)(I\times J)=\mu_n(I)\nu_n(J)
 \longrightarrow\mu(I)\nu(J),
\]
这里用实线分布函数连续点判据。每格取一个标记值，所得有限阶梯和的
积分因格子质量收敛而收敛。对两个测度各用阶梯和近似 $h$ 的积分，
矩形内的两项误差合计不超过 $2\varepsilon$，
矩形外误差不超过 $4\varepsilon\|h\|_\infty$。
先令 $n\to\infty$ 再令 $\varepsilon\downarrow0$，便得联合弱收敛。
\end{solution}

\begin{solution}{probconv-slutsky-studentization}
Slutsky 定理给出 $(T_n,V_n)\xrightarrow d(Z,\sigma^2)$。
在整个 $\mathbb R^2$ 上定义 Borel 函数
\[
 g(t,v)=\begin{cases}t/\sqrt v,&v>0,\\0,&v\le0.\end{cases}
\]
它在所有 $v>0$ 的点连续，可能的不连续点只位于 $v=0$。
极限向量的第二坐标是严格正常数，故极限律不赋予不连续点质量。
连续映射定理于是给出
$g(T_n,V_n)\xrightarrow d g(Z,\sigma^2)=Z/\sigma$。
因 $V_n\ge0$，左侧正是按题意定义的比值。不需要任何独立性；
$\mathbb P(V_n=0)\to0$ 也由 $V_n\xrightarrow{\mathbb P}\sigma^2>0$ 得到。
\end{solution}

\begin{solution}{probconv-vanishing-perturbation}
令 $D_n=Y_n-X_n$。题设给 $D_n\xrightarrow{\mathbb P}0$。
向量版 Slutsky 定理说明
$(X_n,D_n)\xrightarrow d(X,0)$。加法映射
$(x,d)\mapsto x+d$ 连续，因此
$Y_n=X_n+D_n\xrightarrow dX$。在每个 $n$ 内，差 $D_n$ 及联合向量的
定义使用了共同样本空间；不需要 $D_n$ 与 $X_n$ 独立。
\end{solution}

\begin{solution}{probconv-delta-method}
记 $Z_n=r_n(X_n-\theta)$。其分布收敛，所以一致紧。给定
$\varepsilon,\eta>0$，选 $M$ 使
$\sup_n\mathbb P(|Z_n|>M)<\eta$。当 $n$ 大到 $\varepsilon r_n>M$ 时，
\[
 \mathbb P(|X_n-\theta|>\varepsilon)
 =\mathbb P(|Z_n|>\varepsilon r_n)<\eta.
\]
故 $X_n\xrightarrow{\mathbb P}\theta$。

可微性给出
\[
 g(\theta+u)-g(\theta)=g'(\theta)u+u\rho(u),
 \qquad\rho(u)\to0\quad(u\to0),
\]
其中令 $\rho(0)=0$。需要证明 $Z_n\rho(X_n-\theta)$ 依概率趋零。
固定 $\varepsilon,\eta>0$，先选上述 $M>0$，再选 $\delta>0$，使
$|u|\le\delta$ 时 $|\rho(u)|<\varepsilon/M$。则
\[
 \mathbb P(|Z_n\rho(X_n-\theta)|>\varepsilon)
 \le\mathbb P(|Z_n|>M)+\mathbb P(|X_n-\theta|>\delta).
\]
其上极限不超过 $\eta$，再令 $\eta\downarrow0$，得到所需扰动结论。
连续映射给出 $g'(\theta)Z_n\xrightarrow d g'(\theta)Z$，再应用消失扰动
定理，即得结论。即使 $g'(\theta)=0$，证明仍成立，此时一阶尺度的极限
为常数零。
\end{solution}

\begin{solution}{probconv-ui-higher-moment}
由依概率子列原理取 $X_{n_k}\to X$ a.s.，Fatou 引理给出
\[
 \mathbb E|X|^q\le\liminf_k\mathbb E|X_{n_k}|^q
 \le\sup_n\mathbb E|X_n|^q<\infty.
\]
记这个上界为 $C$。在 $|X_n|^p>M$ 上，有
$|X_n|^{p-q}\le M^{-(q-p)/p}$，因此
\[
 \mathbb E\bigl[|X_n|^p\mathbf1_{\{|X_n|^p>M\}}\bigr]
 \le C M^{-(q-p)/p}\longrightarrow0
\]
一致于 $n$。概率空间的 Vitali 判据给出 $L^p$ 收敛。
若取 $X_n=n^{1/q}\mathbf1_{(0,1/n)}$、$X=0$，则依概率收敛且 $q$ 阶
矩恒为一，但没有 $L^q$ 收敛，说明结论不能提升到端点阶数。
\end{solution}

\begin{solution}{probconv-weak-ui-moments}
函数 $x\mapsto|x|^p$ 连续，所以
$|X_n|^p\xrightarrow d|X|^p$。由题设的一致可积性及正文的分布收敛期望
定理，得到
\[
 \mathbb E|X_n|^p\longrightarrow\mathbb E|X|^p<\infty.
\]
这一推理只使用各自的分布，不要求样本空间一致。
即使已有共同空间，也不能据此推出 $L^1$ 收敛。令 $X=Z$、$X_n=-Z$，
其中 $Z$ 为公平符号。则分布始终相同，变量绝对值都为一，故一致可积，
但 $\mathbb E|X_n-X|=2$ 对所有 $n$ 成立。
\end{solution}

\begin{solution}{probconv-ui-scheffe}
因为 $X_n\xrightarrow{\mathbb P}X$，每个子列都有一个 a.s. 收敛再子列。
在这样的再子列上，$\min(X_n,X)\to X$ a.s.，并被可积函数 $X$ 控制，
所以其期望趋于 $\mathbb E X$。数列的子列反证法给出整列
$\mathbb E\min(X_n,X)\to\mathbb E X$。
对非负数的恒等式 $|a-b|=a+b-2\min(a,b)$ 取期望，得到
\[
 \mathbb E|X_n-X|
 =\mathbb E X_n+\mathbb E X-2\mathbb E\min(X_n,X)\longrightarrow0.
\]
Vitali 定理反向部分于是给出 $(X_n)$ 一致可积。非负性使两个均值足以
控制绝对差；一般带符号变量只有均值收敛时不能使用这个论证。
\end{solution}

\begin{solution}{probconv-products-uniform-integrability}
令 $r=(1/p+1/q)^{-1}>1$。Hölder 不等式使用共轭指数 $p/r$、$q/r$，给出
\[
 \mathbb E|X_nY_n|^r
 \le(\mathbb E|X_n|^p)^{r/p}(\mathbb E|Y_n|^q)^{r/q}.
\]
右侧一致有界，因此乘积族有一致的高于一阶的矩，从而一致可积。
两个分量依概率收敛推出联合依概率收敛，再由乘法的连续性，
$X_nY_n\xrightarrow{\mathbb P}XY$。Vitali 定理给出 $XY\in L^1$ 及
乘积的 $L^1$ 收敛。独立性没有被使用，乘积矩的估计来自 Hölder 不等式。
\end{solution}

\begin{solution}{probconv-weak-nonnegative-ui}
对每个固定 $M>0$，$x\mapsto\min\{\max(x,0),M\}$ 是有界连续函数。
非负性及分布收敛给出
\begin{align*}
 \mathbb E(X_n-M)_+
 &=\mathbb E X_n-\mathbb E\min(X_n,M)\\
 &\longrightarrow\mathbb E X-\mathbb E\min(X,M)
 =\mathbb E(X-M)_+.
\end{align*}
给定 $\varepsilon>0$，先选 $M$ 使最后一项小于 $\varepsilon$。
于是存在 $N$，使 $n\ge N$ 时的超额期望小于 $2\varepsilon$。
对有限个可积变量 $X_1,\ldots,X_{N-1}$，再取 $K\ge M$，使它们各自的
$\mathbb E(X_n-K)_+$ 小于 $2\varepsilon$。对 $n\ge N$，增大截断值
也不增加超额期望，故
$\sup_n\mathbb E(X_n-K)_+\le2\varepsilon$。
最后由逐点不等式
\[
 X_n\mathbf1_{\{X_n>2K\}}\le2(X_n-K)_+
\]
得到统一尾积分不超过 $4\varepsilon$，这正是一致可积性。
各步只用单个变量的分布和期望，因此允许不同样本空间。
\end{solution}
